\section{Configuração de treinamento}
\label{sec:treinamento}

Cada experimento é configurado para 10 épocas, com lotes de 32 imagens e
embaralhamento do conjunto de treino a cada época. A função de perda é a
entropia cruzada multiclasse, e o otimizador utilizado é o AdamW, com taxa
de aprendizado constante de $10^{-3}$ e decaimento de pesos de $10^{-2}$.
Os demais hiperparâmetros seguem os valores padrão da implementação:
$\beta_1=0{,}9$, $\beta_2=0{,}999$ e $\epsilon=10^{-8}$. O otimizador recebe
somente os parâmetros definidos como treináveis em cada configuração,
conforme a Tabela~\ref{tab:parametros}.

As execuções utilizam PyTorch 2.14.0 e Torchvision 0.29.0, em uma GPU
NVIDIA GeForce RTX 4050 Laptop GPU, com operações em ponto flutuante de
32 bits. Cada experimento parte dos pesos pré-treinados no ImageNet e
de uma nova camada de classificação. Para favorecer a reprodutibilidade,
é fixada a semente 42 nos geradores aleatórios de Python, NumPy e PyTorch,
e são habilitados algoritmos determinísticos. O embaralhamento das imagens
de treino utiliza um gerador aleatório com semente 42. Os conjuntos de
validação e teste são percorridos em ordem fixa.

Ao final de cada época, o modelo é avaliado no conjunto de validação, e
são registrados a perda e os valores de acurácia top-1 e top-5 de treino e
validação. Os pesos do modelo e as estatísticas do Batch Normalization são
salvos em um arquivo específico para cada época. O conjunto de teste é
avaliado ao término do treinamento, utilizando o modelo da última época.
Durante a validação e o teste, o Dropout e o Stochastic Depth ficam
desativados, e o Batch Normalization utiliza suas estatísticas acumuladas
sem atualizá-las.
